% GENERATED FILE. Edit canonical lessons, then run npm run generate:books.
% canonical-source-hash: fnv1a64:8adb216af56fa6ee
% canonical-lessons: ZH-C530-caipiao, ZH-C530-choujiang, ZH-C530-zhongjiang, ZH-C530-jiangpin, ZH-C530-toujiang

\chapter{Winning a Prize}
\label{ch:zh-winning-a-prize}
\hlchaptermodality{530}

\begin{chapteropening}
\textbf{By the end of this chapter:} \emph{I can talk about lottery tickets, prize draws, winning a prize and trying my luck.}

Complete the last lesson of chapter 530: I can talk about lottery tickets, prize draws, winning a prize and trying my luck.
\end{chapteropening}

\section[cǎipiào]{cǎipiào — lottery ticket}

\label{lesson:ZH-C530-caipiao}



\begin{quote}
Before the new one: say the Mandarin for fortunate, then the Mandarin for fate.
\end{quote}

\subsection*{You'll want to know: cǎipiào}
\emph{cǎipiào} — \textquotedblleft{}lottery ticket\textquotedblright{}.

\begin{quote}
\emph{tā měi ge xīngqī dōu mǎi cǎipiào} — He buys lottery tickets every week.
\end{quote}

Its characters come later; for now it is a word for the ear.

\begin{cousinweb}
\emph{cǎipiào} is built from \emph{cǎi} and \zh{票} 'a ticket'. The first is inherited from Middle Chinese \emph{tshojX} (Guangyun, in Baxter's transcription), glossed 'brilliant colour'; no dictionary we use records an older form of \zh{票}, so its history is not settled.
\end{cousinweb}

\begin{grammarlens}[title={What you've built}]
One more word for saying what would happen.
\end{grammarlens}

\subsection*{Your turn}
\begin{itemize}
\raggedright
  \item \emph{Say it:} \emph{cǎipiào}
  \item \emph{Say it:} \emph{cǎipiào}, once more
  \item \emph{Say it:} \emph{tā měi ge xīngqī dōu mǎi cǎipiào} — He buys lottery tickets every week.
  \item \emph{Recall:} say the Mandarin for fortunate, then the Mandarin for fate, then say \emph{cǎipiào} again
\end{itemize}

\begin{tcolorbox}[breakable,colback=teal!4,colframe=teal!35!black,title={Before you move on}]
What does cǎipiào mean? (\textquotedblleft{}Lottery ticket\textquotedblright{}.) Say it once more.
\end{tcolorbox}
\section[chōujiǎng]{chōujiǎng — to draw for prizes; to hold a prize draw}

\label{lesson:ZH-C530-choujiang}



\begin{quote}
Before the new one: say the Mandarin for fate, then the Mandarin for a lottery ticket.
\end{quote}

\subsection*{You'll want to know: chōujiǎng}
\emph{chōujiǎng} — \textquotedblleft{}to draw for prizes; to hold a prize draw\textquotedblright{}.

\begin{quote}
\emph{zhǐyào mǎi yìbǎi kuài qián de dōngxi jiù kěyǐ chōujiǎng} — Spend a hundred yuan and you can enter the prize draw.
\end{quote}

Its characters come later; for now it is a word for the ear.

\begin{cousinweb}
\emph{chōujiǎng} is built from \emph{chōu} 'to pull out' and \emph{jiǎng}, both inherited from Middle Chinese, \emph{trhjuw} and \emph{tsjangX} (Guangyun, in Baxter's transcription), which glosses the second 'to encourage' and 'to praise'; the sense 'a prize' is not given there.
\end{cousinweb}

\begin{grammarlens}[title={What you've built}]
One more word for saying what would happen.
\end{grammarlens}

\subsection*{Your turn}
\begin{itemize}
\raggedright
  \item \emph{Say it:} \emph{chōujiǎng}
  \item \emph{Say it:} \emph{chōujiǎng}, once more
  \item \emph{Say it:} \emph{zhǐyào mǎi yìbǎi kuài qián de dōngxi jiù kěyǐ chōujiǎng} — Spend a hundred yuan and you can enter the prize draw.
  \item \emph{Recall:} say the Mandarin for fate, then the Mandarin for a lottery ticket, then say \emph{chōujiǎng} again
\end{itemize}

\begin{tcolorbox}[breakable,colback=teal!4,colframe=teal!35!black,title={Before you move on}]
What does chōujiǎng mean? (\textquotedblleft{}To draw for prizes; to hold a prize draw\textquotedblright{}.) Say it once more.
\end{tcolorbox}
\section[zhòngjiǎng]{zhòngjiǎng — to win in a draw; to win a prize}

\label{lesson:ZH-C530-zhongjiang}



\begin{quote}
Before the new one: say the Mandarin for a lottery ticket, then the Mandarin for to draw for prizes.
\end{quote}

\subsection*{You'll want to know: zhòngjiǎng}
\emph{zhòngjiǎng} — \textquotedblleft{}to win in a draw; to win a prize\textquotedblright{}.

\begin{quote}
\emph{yàoshi wǒ zhòngjiǎng le wǒ jiù qù lǚyóu} — If I win a prize, I'll go travelling.
\end{quote}

Its characters come later; for now it is a word for the ear.

\begin{cousinweb}
\emph{zhòngjiǎng} is built from \zh{中} zhòng 'to hit the mark' and \emph{jiǎng}, both inherited from Middle Chinese, \emph{trjuwngH} and \emph{tsjangX} (Guangyun, in Baxter's transcription), which glosses the first 'to hit' and the second 'to encourage'.
\end{cousinweb}

\begin{grammarlens}[title={What you've built}]
One more word for saying what would happen.
\end{grammarlens}

\subsection*{Your turn}
\begin{itemize}
\raggedright
  \item \emph{Say it:} \emph{zhòngjiǎng}
  \item \emph{Say it:} \emph{zhòngjiǎng}, once more
  \item \emph{Say it:} \emph{yàoshi wǒ zhòngjiǎng le wǒ jiù qù lǚyóu} — If I win a prize, I'll go travelling.
  \item \emph{Recall:} say the Mandarin for a lottery ticket, then the Mandarin for to draw for prizes, then say \emph{zhòngjiǎng} again
\end{itemize}

\begin{tcolorbox}[breakable,colback=teal!4,colframe=teal!35!black,title={Before you move on}]
What does zhòngjiǎng mean? (\textquotedblleft{}To win in a draw; to win a prize\textquotedblright{}.) Say it once more.
\end{tcolorbox}
\section[jiǎngpǐn]{jiǎngpǐn — a prize you win; an award}

\label{lesson:ZH-C530-jiangpin}



\begin{quote}
Before the new one: say the Mandarin for to draw for prizes, then the Mandarin for to win in a draw.
\end{quote}

\subsection*{You'll want to know: jiǎngpǐn}
\emph{jiǎngpǐn} — \textquotedblleft{}a prize you win; an award\textquotedblright{}.

\begin{quote}
\emph{rúguǒ nǐ zhòngjiǎng le jiù kěyǐ ná dào jiǎngpǐn} — If you win, you'll get a prize.
\end{quote}

Its characters come later; for now it is a word for the ear.

Two third tones meet in this word: jiǎng is said with a rising tone before pǐn, and the spelling keeps both marks.

\begin{cousinweb}
\emph{jiǎngpǐn} is built from \emph{jiǎng} and \emph{pǐn}, both inherited from Middle Chinese, \emph{tsjangX} and \emph{phimX} (Guangyun, in Baxter's transcription), which glosses the first 'to encourage' and the second 'a rank' and 'a kind'.
\end{cousinweb}

\begin{grammarlens}[title={What you've built}]
One more word for saying what would happen.
\end{grammarlens}

\subsection*{Your turn}
\begin{itemize}
\raggedright
  \item \emph{Say it:} \emph{jiǎngpǐn}
  \item \emph{Say it:} \emph{jiǎngpǐn}, once more
  \item \emph{Say it:} \emph{rúguǒ nǐ zhòngjiǎng le jiù kěyǐ ná dào jiǎngpǐn} — If you win, you'll get a prize.
  \item \emph{Recall:} say the Mandarin for to draw for prizes, then the Mandarin for to win in a draw, then say \emph{jiǎngpǐn} again
\end{itemize}

\begin{tcolorbox}[breakable,colback=teal!4,colframe=teal!35!black,title={Before you move on}]
What does jiǎngpǐn mean? (\textquotedblleft{}A prize you win; an award\textquotedblright{}.) Say it once more.
\end{tcolorbox}
\section[tóujiǎng]{tóujiǎng — first prize; the jackpot}

\label{lesson:ZH-C530-toujiang}



\begin{quote}
Before the new one: say the Mandarin for to win in a draw, then the Mandarin for a prize you win.
\end{quote}

\subsection*{You'll want to know: tóujiǎng}
\emph{tóujiǎng} — \textquotedblleft{}first prize; the jackpot\textquotedblright{}.

\begin{quote}
\emph{zhè cì chōujiǎng de tóujiǎng shì yìbǎi wàn} — The first prize in this draw is a million.
\end{quote}

Its characters come later; for now it is a word for the ear.

\begin{cousinweb}
\emph{tóujiǎng} is built from \zh{头} 'the head' and \emph{jiǎng}, both inherited from Middle Chinese, \emph{duw} and \emph{tsjangX} (Guangyun, in Baxter's transcription), which glosses the second 'to encourage': the reward at the head of the list.
\end{cousinweb}

\begin{grammarlens}[title={What you've built}]
One more word for saying what would happen.
\end{grammarlens}

\subsection*{Your turn}
\begin{itemize}
\raggedright
  \item \emph{Say it:} \emph{tóujiǎng}
  \item \emph{Say it:} \emph{tóujiǎng}, once more
  \item \emph{Say it:} \emph{zhè cì chōujiǎng de tóujiǎng shì yìbǎi wàn} — The first prize in this draw is a million.
  \item \emph{Recall:} say the Mandarin for to win in a draw, then the Mandarin for a prize you win, then say \emph{tóujiǎng} again
\end{itemize}

\begin{tcolorbox}[breakable,colback=teal!4,colframe=teal!35!black,title={Before you move on}]
What does tóujiǎng mean? (\textquotedblleft{}First prize; the jackpot\textquotedblright{}.) Say it once more.
\end{tcolorbox}
